%!TEX root = ../template.tex %
% !TeX spellcheck = en_US
%%%%%%%%%%%%%%%%%%%%%%%%%%%%%%%%%%%%%%%%%%%%%%%%%%%%%%%%%%%%%%%%%%%%
%% acknowledgements.tex
%% NOVA thesis document file
%%
%% Text with acknowledgements
%%%%%%%%%%%%%%%%%%%%%%%%%%%%%%%%%%%%%%%%%%%%%%%%%%%%%%%%%%%%%%%%%%%%

\typeout{NT FILE acknowledgements.tex}%

\begin{ntacknowledgements}

This PhD thesis would not have been possible without the valuable guidance from my supervisors Prof.~Christos Tserkezis, Prof.~Nuno~M.~R.~Peres and Prof.~N.~Asger Mortensen.
I am deeply grateful for them to accept me into POLIMA and support me throughout the PhD studies.
I am extremely thankful to Christos for his daily guidance and availability, his feedback given on every occasion and everytime I needed, the time he spent in proofreading this thesis,
and for teaching me so much about hyphenization.
I want to express my gratitude towards Nuno, who brought me into a new field in physics I was not familiar with, and by believing in my capacities.
I feel lucky for receiving his technical guidance, advice, and wisdom that I feel like would not be possible to receive in other circumstances.
I am deeply thankful as well to Asger as well, who trusted in my worth and accepted me for my PhD studies when no one else did, and for teaching me how to appreciate whine.
That is very dear to me.

I would like to thank my colleagues in POLIMA, who provided a great social and scientific environment to thrive at SDU during the past three years.
Especially those that helped me by proofreading the thesis, and Louise for her unwavering support during the roughest times.

A big part of this thesis would also not be possible without the guidance and advice from Prof.~Juan José Palacios, who received me in his group in Madrid,
 whom I had the opportunity to learn a lot from, and for his patience too.
I also appreciate the risk he took in accepting someone from a group he had never collaborated with before.
I was very well received in Madrid with the Spanish warmth, and loved to engage with everyone UAM\@.
I would like to give credit to Alex for his work that set the staged for a big part of this thesis, to thank Maurício for testing the code, and for Guilherme to provide me with the Wannier models that were so useful for the work
I started in Madrid.

I want also to thank the members of my PhD evaluation committee, Prof.~Kristian Sommer Thygesen and Prof.~Diego Martín Cano for taking some of their time to assess this work and for reading my thesis.

Less professionally, I would like to say that the 8th floor at Campus Kollegiet made my life so much easier during hard times, and thank them for making my life much more interesting socially.
A special thanks to Nida, who I trusted in so much.
A quick thanks to Andreas, Eja, Maria, Bianca, Shneha, Eda, Victor, Malene, Luigi, Sole, Cecilia, Jenistan, and Christoffer.
They all made the experience of living in the 8th very funny.
I apologize if I forgot anyone.

On a personal level, I want to say that the years spent at university with MEFT at IST are among the best years of my life.
I am very grateful to all the friends I made along the way, who made it easier to finish the degree.
No one makes MEFT alone.

I want to extend my thanks to my high-school friends, who have been around for so long.
I hold their friendship really dearly.
I also want to express my gratitude towards my high-school physics teacher, who inspired me to pursue physics in university.

Finally, I want to acknowledge the people in my life that have always been there, in the good and bad times, my family.
For their unconditional love and support, and for helping me during these hard times abroad, even with them being in the far field.


The author acknowledges support by the Otto M{\o}nsted Foundation (grant No. 24-70-1631) and the Independent Research Fund Denmark (grant No. 2032-00045A).
%
The Center for Polariton-driven Light--Matter Interactions (POLIMA) is funded by the Danish National Research Foundation (project No.~DNRF165).


\end{ntacknowledgements}